\documentclass[a4paper,12pt]{article}
\usepackage[utf8]{inputenc}
\usepackage[ngerman]{babel}
\usepackage{amsmath}
\usepackage{geometry}
\geometry{a4paper, top=15mm, left=15mm, right=15mm, bottom=15mm,
	headsep=10mm, footskip=12mm}
\usepackage{graphicx}
\usepackage{booktabs}
\usepackage{parskip}
\usepackage{href-ul}
\hypersetup{
	colorlinks=true,
	linkcolor=blue,
	urlcolor=blue}
\title{Physik mit Humor: Gewichtskräfte überall!}
\date{}

\begin{document}
	
	
	\section*{Formel des Tages}
	
	\[
	F = m \cdot g
	\]
	\begin{itemize}
		\item \( F \): Gewichtskraft in Newton (N)
		\item \( m \): Masse in Kilogramm (kg)
		\item \( g \): Ortsfaktor / Fallbeschleunigung in \( \frac{\text{m}}{\text{s}^2} \)
	\end{itemize}
	
	\hrulefill
	
	\section{1. Der hungrige Hamster auf dem Mond }
	
  Hamster Henry wiegt 150 g. Auf dem Mond ist die Schwerkraft nur \( g = 1{,}6\, \frac{\text{m}}{\text{s}^2} \).
	
	\begin{itemize}
		\item a) Berechne Henrys Gewichtskraft auf der Erde (\( g = 9{,}81\, \frac{\text{m}}{\text{s}^2} \)):
		
		\vspace{1em}
		\( m = \underline{\hspace{4cm}} \)
		
		\( F_\text{Erde} = \underline{\hspace{6cm}} \)
		
		\item b) Berechne Henrys Gewichtskraft auf dem Mond:
		
		\vspace{1em}
		\( F_\text{Mond} = \underline{\hspace{6cm}} \)
		
		\item c) Um wie viel Prozent ist Henry dort „leichter“?
		
		\vspace{1em}
		\( \text{Differenz in \%} = \underline{\hspace{6cm}} \)
	\end{itemize}
	
	\hrulefill
	
	\section*{2. Die Yoga-Einlage des Elefanten }
	
	Ein Elefant (Masse 3000 kg) steht bei seiner Yoga-Übung auf einem Bein auf einer Waage.
	
	\begin{itemize}
		\item a) Welche Kraft wirkt auf das eine Bein?  
		\( F = \underline{\hspace{6cm}} \)
		
		\item b) die Waage zeigt 33 kN an. Berechne die zugehörige Masse (mit \( g = 10 \, \frac{\text{m}}{\text{s}^2} \)):  
		\( m = \underline{\hspace{6cm}} \)
	\end{itemize}
	
	\hrulefill
	
	\section*{3. Der Heliumballon in der Sauna }
	
	Ein Ballon mit einer Masse von 50 g liegt in der Sauna.
	
	\begin{itemize}
		\item a) Berechne seine Gewichtskraft auf der Erde:
		
		\vspace{1em}
		\( m = \underline{\hspace{4cm}} \)
		
		\( F = \underline{\hspace{6cm}} \)
	\end{itemize}
	
	\hrulefill
	
	\section*{4. Die intergalaktische Kofferwaage }
	
	Ein Koffer (Masse 25 kg) wird auf verschiedenen Planeten gewogen.
	
	\begin{center}
		\begin{tabular}{l c}
			\toprule
			\textbf{Planet} & \textbf{g in } \( \frac{\text{m}}{\text{s}^2} \) \\
			\midrule
			Erde & 9,81 \\
			Mars & 3,71 \\
			Jupiter & 24,79 \\
			Uranus & 8,69 \\
			\bottomrule
		\end{tabular}
	\end{center}
	
	\begin{itemize}
		\item a) Trage für jeden Planeten die Gewichtskraft in N ein:
		
		\vspace{0.5em}
		\textbf{Erde:} \( F = \underline{\hspace{6cm}} \)\\
		\textbf{Mars:} \( F = \underline{\hspace{6cm}} \)\\
		\textbf{Jupiter:} \( F = \underline{\hspace{6cm}} \)\\
		\textbf{Uranus:} \( F = \underline{\hspace{6cm}} \)
		
		\item b) Wo ist der Koffer am schwersten? \underline{\hspace{4cm}}
		 
		\underline{\hspace{10cm}}
	\end{itemize}
	
	\hrulefill
	
	\section*{5. Das schwer verliebte Hufeisen }
	
	Ein Hufeisen mit Masse 2 kg hängt an einem Magneten und zieht ein Eisenstück nach oben.
	
	\begin{itemize}
		\item a) Berechne die Gewichtskraft des Hufeisens:
		
		\( F = \underline{\hspace{6cm}} \)
		
		\item b) Welche Mindestkraft muss der Magnet aufbringen?  
		\( F_\text{Magnet} = \underline{\hspace{6cm}} \)
		
		\item c) Philosophieren erlaubt:  
		Warum ist Liebe manchmal messbar?
		
		\vspace{2em}
		\underline{\hspace{12cm}}  
	\end{itemize}
	
	\fbox{
		\begin{minipage}{0.5\textwidth}
			Zur Lösung bitte \href{https://www.okuyakl.de/phys/p7fmgL058/ll058.pdf}{hier klicken} oder den QR-Code scannen.\\
			Weitere Arbeitsblätter gibt es unter 
			
			\href{https://www.okuyakl.de}{www.okuyakl.de}
		\end{minipage}
		\hfill
		\begin{minipage}{0.4\textwidth}
			\includegraphics[width=1.5 cm]{../../viecher/zwe03}
			\includegraphics[width=3 cm]{qr058}
			\includegraphics[width=2 cm]{../../viecher/afanticon1}
			
	    \end{minipage}}

	\end{document}
	
\end{document}
